\section[Pfister-Roumy]{Création de la liste Pfister-Roumy}

La  bibliothèque Cujas fut un \enquote{Centre d'acquisition et de diffusion de l'information scientifique et technique} (CADIST). Habituellement, les bibliothèques possèdent des fonds patrimoniaux issu pour certaines de confiscations révolutionnaires. Hors, pour Cujas, sa collection d'origine \enquote{de la Faculté de Droit de Paris a été dispersée sous la Révolution}. Le fonds actuel a été reconstruit \enquote{rétropectivement par des bibliothécaires, dans le but de reconstituer \enquote{le cabinet de travail d'un robin d'Ancien Régime} - telle était, du moins, la doxa transmise jusqu'à présent de bouche à oreille de bibliothécaire}. Le fonds patrimonial de Cujas s'est constitué grâce à des achats, dons et legs, essentiellement à partir des années 1870. La constitution de ce fonds comprend : \enquote{12 000 volumes imprimés, datant de la fin du XVe siècle au début du XIXe siècle, essentiellement des livres des XVIIe et XVIIIe siècles, 7 incunables, quelques ouvrages précieux postérieurs à 1800 et quelques périodiques. Le quart de ce fonds provient de pays étrangers : Allemagne, Pays-Bas, Angleterre et Italie. Particulièrement riche, il comprend 25 \% d'ouvrages non référencés à la BnF. La réserve conserve également des cours et notes de cours manuscrits, ainsi qu'une partie des archives de la Faculté de Droit.}. Ces fonds patrimoniaux font en sorte de représenter les grands courants du droit : droit civil, criminel et canonique, pratique du droit ancien ou nouveau (après 1789), coutumiers, oeuvres des grands jurisconsultes français et étrangers.
	
Pour mettre en avant la richesse de ce patrimoine, Cujas a l'idée de créer une bibliothèque numérique. Cette idée germe au moment de la numérisation des cours de Jean Carbonnier pour l'exposition virtuelle. L'objectif de la bibliothèque Cujas est d'utiliser le scanner qu'elle venait d'acheter, au moment de sa candidature au projet Persée. A ce moment-là, Persée offrait la possibilité au candidat retenu de devenir un pôle de réplication. Malheureusement, le projet n'a pas abouti. Toutefois, la bibliothèque Cujas souhaitant utiliser le scanner qu'elle avait acquis prend la décision de s'associer avec deux professeurs de la Faculté de droit de Paris, Laurent Pfister et Franck Roumy. Cette collaboration permet de dresser la première liste de référence, intitulé plus tard la liste Pfister-Roumy, visant à numériser en priorité les cours d'histoire du droit.

L'objectif de la liste Pfister-Roumy est de numériser uniquement des monographies ou des périodiques provenant \enquote{uniquement des collections de Cujas}. Cette liste porte sur les \enquote{fondamentaux du droit à numériser en priorité}. Les deux professeurs en ont compté à l'origine \enquote{471 titres fondamentaux pour la recherche en histoire du droit}. Ces ouvrages couvrent la période du \enquote{XVe siècle au XXe siècle, soit 1462 volumes et 1 001 559 pages}. A cette liste, ils ajoutent deux périodiques : \enquote{Duvergier et le Bulletin des arrêts de la Cour de cassation (Chambre criminelle et Chambres civiles) jusqu'en 1950.}. Il s'agit du projet d'origine penser au moment de la candidature de Cujas pour devenir pôle de réplication de Persée. Ce projet se veut ambitieux, à travers l'utilisation d'une \enquote{numérisation automatique non descriptive, en mode image océrisé, avec une recherche plein texte et la structuration documentaire à granularité fine, emploi de la DTD TEI, réalisé exclussivement en interne}. 

L'abandon de ce projet oblige la bibliothèque Cujas a repensé ce projet de numérisation, sauf que le nombre de volumes numérisés est revu à la baisse. Celui-ci passe de 471 ouvrages à 436 ouvrages. En effet, à la même époque, la BnF réalise un projet similaire à celui de Cujas en parallèle, intitulé la liste Kerbrat. Cette liste est formé par le professeur Yann Kerbrat et d'un groupe de quatorze universitaire avec comme objectif \enquote{de créer une liste de références à partir des collections de la BnF} sur le droit. Dans un soucis de moyen financier et de pertinence, des choix sont fait entre les deux institutions afin de décider qui numérisé quoi. Cela signifie que lorsque la bibliothèque Cujas possède un ouvrage identique à celui de la BnF, pour éviter les incohérences chronologiques, il a été décidé de leur permettre de numériser l'ouvrage aussi mais dans une édition différente de celle conservée à la BnF. D'autres ouvrages ont simplement été retirer ne pouvant pas être remplacer par une autre édition. 

Une autre contrainte vient du scanner en lui-même utilisé par la bibliothèque Cujas qui s'avère moins efficace les amenant à rencontrer de \enquote{nombreux problèmes liés à son utilisation : il n'est pas possible d'utiliser toutes les fonctions automatiques de numérisation du scanner, forte probabilité de devoir passer en manuel pour les papiers anciens, difficultés de réglage et fréquents dysfonctionnements nécessitant d'éteindre et de rallumer les appareils et de refaire tous les réglages à chaques fois}. Le scanner présent aussi un fort risque de dégradation des documents originaux au moment de la numérisation. C'est pour cela que lorsque c'est possible elle délègue la numérisation des volumes de sa liste. Elle a l'objectif de créer une bibliothèque numérique juridique fédérée entre les documents de Cujas et de ses bibliothèques partenaires à travers la mise en place d'un moissonnage OAI. Malgré les difficultés présenties face à la taille de ce projet, la bibliothèque Cujas les a pris en compte et a validé sa liste Pfister-Roumy. Une fois la validation obtenue, la bibliothèque utilise sa chaîne de numérisation interne qu'elle avait mis en place pour l'exposition de Jean Carbonnier et entreprend la numérisation de sa liste tout en déléguant une partie à un prestataire externe. Cette liste Pfister-Roumy représente l'un des corpus les plus importants qui sera à l'origine de la création de la bibliothèque numérique de Cujas, CujasNum.

